\chapter{Introducción}

La transformación digital del Estado peruano ha avanzado de manera desigual entre
sus distintas funciones: mientras que los procesos transaccionales y administrativos
han incorporado herramientas de gestión digital, los procesos de análisis técnico-legal
de alta complejidad ---como la revisión de contratos de concesión de infraestructura---
permanecen fundamentalmente dependientes del criterio humano y de métodos manuales
no estandarizados. Esta brecha entre la complejidad creciente de los instrumentos
jurídicos y la capacidad de análisis disponible representa un riesgo latente para
la calidad de las decisiones y la eficiencia del ciclo de inversión pública.

A nivel global, la investigación en \textit{Natural Language Processing} (NLP)
aplicado al dominio legal ha producido avances significativos: conjuntos de datos
especializados como CUAD \cite{hendrycks2021cuad} y ContractNLI
\cite{koreeda2021contractnli}, plataformas comerciales de análisis documental
\cite{google_contract_docai_2021, ibm_watsonx_contracts_2024}, y revisiones
sistemáticas sobre el estado del arte en LLMs para derecho \cite{singh2025survey}.
Sin embargo, estas soluciones presentan una limitación común frente al contexto
peruano: están diseñadas para marcos normativos anglosajones, no incorporan la
normativa sectorial local (Ley de APP, TUO de Contrataciones del Estado, normas
del OSITRAN, OSINERGMIN, SUNASS, entre otras), y no están adaptadas a los
patrones estructurales específicos de los contratos de concesión elaborados por
instituciones como ProInversión. Esta brecha de adaptación motivó el presente proyecto.

\textit{ContractIA} surge como respuesta a esa brecha: un sistema de inteligencia
artificial diseñado desde el inicio para el dominio legal peruano, con foco en la
auditoría automatizada de contratos de concesión de infraestructura. El proyecto
se enmarca en el curso de Capstone Project II de la Maestría y constituye la
evidencia integradora de competencias en inteligencia artificial aplicada,
arquitectura de sistemas de datos y gestión de productos digitales. Su propósito
no es reemplazar al especialista legal, sino dotarlo de una herramienta de
\textit{screening} sistemático que eleve la cobertura, la consistencia y la
trazabilidad del proceso de revisión, reduciendo el tiempo invertido en tareas
de detección para concentrar el juicio experto en la interpretación y decisión.

El presente informe documenta el recorrido completo desde la identificación del
problema hasta la validación del MVP en producción, organizado en las siguientes
secciones: la Sección~2 desarrolla el contexto organizacional del problema y las
entrevistas con especialistas; la Sección~3 define los objetivos SMART del proyecto;
la Sección~4 describe el marco metodológico adoptado; la Sección~5 presenta el
autodiagnóstico del proceso actual; la Sección~6 define el MVP y sus criterios de
aceptación; la Sección~7 detalla la arquitectura técnica de la solución; la
Sección~8 describe el proceso de construcción e implementación; la Sección~9
establece la métrica principal de validación; la Sección~10 reporta los resultados
de la evaluación experimental; la Sección~11 analiza las limitaciones y riesgos
identificados; la Sección~12 sistematiza los aprendizajes del proyecto; la
Sección~13 presenta las conclusiones; la Sección~14 proyecta el trabajo futuro
mediante un roadmap de cuatro fases; y las Secciones~15 y~16 consolidan las
referencias y los anexos técnicos.